\chapter{Ejemplos OpenAPI y BIAN}
\label{anexo:f}

Este anexo fija el formato de un par y los scores del piloto de siete instancias. El extracto
de \textit{Payment Initiation} muestra cómo se lee un contrato. No reemplaza a I1--I6.

\section{Ejemplo reducido (unidad de análisis)}

\begin{table}[H]
  \centering
  \caption{Ejemplo reducido — pares OpenAPI $\leftrightarrow$ BIAN}
  \label{tab:anexo-f-reducido}
  \small
  \begin{tabular}{@{}p{0.44\textwidth}p{0.44\textwidth}@{}}
    \toprule
    \textbf{Elemento OpenAPI} & \textbf{Elemento BIAN (referencia)} \\
    \midrule
    \texttt{POST /payment-orders} & \Behavior{} \texttt{Initiate} \\
    \texttt{GET /payment-orders/\{id\}} & \Behavior{} \texttt{Retrieve} \\
    Esquema \texttt{PaymentOrder} (\texttt{amount}, \texttt{currency}) &
    \BusinessObject{} \texttt{Payment} / \texttt{PaymentAmount} \\
    \bottomrule
  \end{tabular}
\end{table}

\section{Ejemplo ampliado (fragmento de contrato)}

\subsection{Contrato OpenAPI (extracto ilustrativo)}

\lstinputlisting{../datos-objeto-estudio/ejemplo_openapi_payment_initiation.yaml}

\subsection{Referencia BIAN (Payment Initiation — ilustrativo)}

\begin{table}[H]
  \centering
  \caption{Elementos BIAN emparejables con el extracto OpenAPI}
  \label{tab:anexo-f-bian}
  \small
  \begin{tabular}{@{}p{0.28\textwidth}p{0.62\textwidth}@{}}
    \toprule
    \textbf{Tipo BIAN} & \textbf{Elemento de referencia} \\
    \midrule
    \ServiceDomain{} & Payment Initiation (dominio de pagos) \\
    \Behavior{} & \texttt{Initiate}, \texttt{Retrieve}, \texttt{Update} \\
    \BusinessObject{} & \texttt{PaymentInstruction}, \texttt{PaymentAmount}, \texttt{Party} \\
    Relación de control & Orquestación iniciación $\rightarrow$ confirmación de estado \\
    \bottomrule
  \end{tabular}
\end{table}

\subsection{Correspondencias para el modelo S/E/C}

\begin{table}[H]
  \centering
  \caption{Lectura del ejemplo bajo S, E y C}
  \label{tab:anexo-f-scores}
  \small
  \begin{tabular}{@{}p{0.12\textwidth}p{0.78\textwidth}@{}}
    \toprule
    \textbf{Score} & \textbf{Qué evaluaría en este extracto} \\
    \midrule
    E &
    Correspondencia \texttt{/payment-initiations} $\leftrightarrow$ \Behavior{};
    anidamiento \texttt{components/schemas} $\leftrightarrow$ \BusinessObject{} \\
    S &
    Similitud \texttt{instructedAmount/currency} $\leftrightarrow$ \texttt{PaymentAmount};
    \texttt{creditor} $\leftrightarrow$ \texttt{Party} \\
    C &
    Proporción de \Behavior{} y \BusinessObject{} del SD cubiertos por operaciones y esquemas
    del fragmento \\
    \bottomrule
  \end{tabular}
\end{table}

\section{Scores del piloto (siete instancias)}

Pesos iguales, $\alpha=\beta=\gamma=\frac{1}{3}$. S es el coseno MiniLM sobre los pares que
congeló E. El S léxico (Jaccard + TF-IDF) quedó entre 0,00 y 0,03 en I1--I5, y en 0 en I6.
Los umbrales 0,80 / 0,60 / 0,40 no se recalibraron: son la convención de §1.4. Nadie llega a
Alta.

\begin{table}[H]
  \centering
  \caption{Piloto I1--I6 --- S (MiniLM), E, C y \AlignmentScore{}}
  \label{tab:anexo-f-escenarios}
  \footnotesize
  \begin{tabular}{@{}lp{4.4cm}ccccc@{}}
    \toprule
    \textbf{ID} & \textbf{Par} & \textbf{S} & \textbf{E} & \textbf{C} & \textbf{Score} & \textbf{Nivel} \\
    \midrule
    I1 & accounts 2.4.2 $\times$ Savings Account & 0,56 & 0,35 & 0,40 & 0,44 & Baja \\
    I2 & accounts 2.4.2 $\times$ Current Account & 0,53 & 0,47 & 0,50 & 0,50 & Baja \\
    I3a & opendata-accounts $\times$ Savings Account & 0,45 & 0,56 & 0,20 & 0,41 & Baja \\
    I3b & opendata-accounts $\times$ Current Account & 0,46 & 0,56 & 0,20 & 0,41 & Baja \\
    I4 & payments 4.0.0 $\times$ Payment Order & 0,52 & 0,29 & 0,25 & 0,35 & Nula \\
    I5 & automatic-payments $\times$ Payment Order & 0,59 & 0,32 & 0,50 & 0,47 & Baja \\
    I6 & enrollments $\times$ Payment Order & 0,00 & 0,00 & 0,00 & 0,00 & Nula \\
    \bottomrule
  \end{tabular}
\end{table}

I6 en cero es el control negativo: no hubo par. Precision, recall y F1 no están en esta tabla.
Salen del subconjunto etiquetado (§4.2), que todavía no está cerrado.

\section{Nota metodológica}

El procedimiento de normalización, matching y scores está en el Capítulo~4. Este anexo deja el
ejemplo de lectura y los siete números del piloto.
